% Vertical profiles are conceptual normalized shapes, not AFGL measurements.
\begin{figure}[H]
\centering
\begin{tikzpicture}[x=1cm,y=1cm,font=\sffamily\fontsize{9}{11}\selectfont,>=Latex]
\node[align=center,font=\sffamily\bfseries\fontsize{10}{12}\selectfont] at (2.4,5.7) {\FigLang{1. Rayleigh 산란의 크기}{1. Rayleigh scattering amount}};
\node[align=center,font=\sffamily\bfseries\fontsize{10}{12}\selectfont] at (7.65,5.7) {\FigLang{2. 기체의 수직 분포}{2. Vertical gas distribution}};
\node[align=center,font=\sffamily\bfseries\fontsize{10}{12}\selectfont] at (12.9,5.7) {\FigLang{3. 기체의 총량}{3. Total gas column}};
\foreach \x in {0.6,2.8}{\fill[ocrtTeal!6] (\x,1.8) rectangle +(1.2,3.1);\draw[ocrtLine] (\x,1.8) rectangle +(1.2,3.1);}
\foreach \px/\py in {0.83/2.1,1.4/2.2,1.05/2.6,1.5/2.9,0.9/3.25,1.4/3.6,1.1/4.1,1.5/4.45}{\fill[ocrtTeal] (\px,\py) circle (2.4pt);}
\foreach \px/\py in {3.08/2.1,3.58/2.65,3.2/3.35,3.56/4.05}{\fill[ocrtTeal] (\px,\py) circle (2.4pt);}
\node[align=center] at (1.2,1.38) {$P_0$};
\node[align=center] at (3.4,1.38) {$P_0/2$};
\draw[->,ocrtTeal,thick] (1.7,4.9)--(2.45,4.9);
\node[align=center,text width=4.6cm] at (2.4,0.43) {$\tau_R(\lambda,P)=\tau_R(\lambda,P_0)\dfrac{P}{P_0}$\\$P_0=1013.25\;\mathrm{hPa}$};
% Two alternative normalized shape sketches; equal column is merely illustrative.
\draw[->,ocrtGray] (5.8,1.8)--(9.8,1.8) node[below left] {$g(z)$};
\draw[->,ocrtGray] (5.8,1.8)--(5.8,5.0) node[above] {$z$};
\draw[ocrtTeal,very thick] plot[smooth] coordinates {(8.9,1.8)(8.0,2.4)(7.2,3.0)(6.6,3.6)(6.1,4.6)};
\draw[orange!85!black,very thick,dashed] plot[smooth] coordinates {(7.0,1.8)(7.2,2.4)(7.7,3.0)(7.55,3.6)(6.2,4.6)};
\node[align=center,text width=4.75cm] at (7.65,1.04) {\FigLang{프로파일 선택 = 형태 선택}{Profile selection = shape choice}};
\node[align=center,text width=4.7cm] at (7.65,0.35) {\texttt{--atm-profile}\\\FigLang{예: \texttt{usstd76}, \texttt{tropical}}{e.g. \texttt{usstd76}, \texttt{tropical}}};
% Scaling the same shape preserves height-dependent relative distribution.
\draw[->,ocrtGray] (11.1,1.8)--(15.3,1.8) node[below left] {$n(z)$};
\draw[->,ocrtGray] (11.1,1.8)--(11.1,5.0) node[above] {$z$};
\fill[ocrtTeal!9] (11.1,1.8)--(14.7,1.8)--(13.9,2.4)--(13.1,3.0)--(12.5,3.6)--(11.5,4.6)--(11.1,4.6)--cycle;
\draw[ocrtTeal,very thick] plot[smooth] coordinates {(12.9,1.8)(12.5,2.4)(12.1,3.0)(11.8,3.6)(11.3,4.6)};
\draw[orange!85!black,very thick,dashed] plot[smooth] coordinates {(14.7,1.8)(13.9,2.4)(13.1,3.0)(12.5,3.6)(11.5,4.6)};
\draw[->,ocrtGray,thick] (12.5,2.42)--(13.8,2.42) node[midway,above] {$\times k$};
\node[align=center,text width=4.8cm] at (12.9,1.02) {$N=\int n(z)\,dz$};
\node[align=center,text width=4.85cm] at (12.9,0.28) {\texttt{--gas-column-*}\\\FigLang{형태를 보존하며 총량 조절}{Scale column; preserve shape}};
\end{tikzpicture}
\caption{\FigLang{\texttt{--pressure}는 Rayleigh 광학두께를, \texttt{--atm-profile}은 기체흡수 프로파일을, \texttt{--gas-column-*}는 각 기체의 총량을 조절한다. 그림의 점·곡선은 개념도이며 실제 AFGL 값이 아니다. 압력 0만으로 기체흡수는 꺼지지 않는다. 컬럼 단위는 H$_2$O: g cm$^{-2}$, O$_3$/NO$_2$: DU, O$_2$/CO$_2$/CH$_4$: molecules cm$^{-2}$다.}{\texttt{--pressure} scales Rayleigh optical depth, \texttt{--atm-profile} selects the gas-absorption profile, and \texttt{--gas-column-*} sets each gas column. Dots and curves are conceptual, not AFGL data. Pressure 0 does not disable gas absorption. Column units are g cm$^{-2}$ for H$_2$O, DU for O$_3$/NO$_2$, and molecules cm$^{-2}$ for O$_2$/CO$_2$/CH$_4$.}}
\label{fig:options-gas}
\end{figure}
